\documentclass{beamer}

\usetheme{Warsaw}

\usepackage{mathlist}

\usepackage[frenchb]{babel}
\usepackage[utf8]{inputenc}
\usepackage{amsmath,amssymb}

\usepackage{pstricks}

\usepackage{times}

\usepackage{ulem}

\def\bxi{\boldsymbol\xi}
\def\bomega{\boldsymbol\omega}

\title{Exercices - Correction}

\author[Fabian Bastin]{Fabian Bastin\\\url{fabian.bastin@cirrelt.ca}\\Université de Montréal -- CIRRELT}

\date{ENAC -- Hiver 2015}

\setbeamertemplate{footline}[frame number]

\begin{document}

\frame{\titlepage}

\begin{frame}
\frametitle{Exercice 1}

%(Birge et Louveaux, page 87)
Considérons le programme de seconde étape
\[
Q(x, \xi) = \min_y \lbrace y \ |\ \xi y = 1-x, y \geq 0 \rbrace.
\]
On suppose que $\bxi$ suit une distribution triangulaire sur $[0,1]$,
à savoir $P[\bxi \leq u] = u^2$.

\mbox{}

{\blue (a)}
Peut-on parler de recours fixe? Pourquoi?

\mbox{}

On ne peut parler de recours fixe puisqu'ici, en suivant
les notations adoptées en cours, $W \equiv \bxi$, et donc est aléatoire.

Comme $\xi$ peut prendre la valeur 0, la transformation
\[
y = 1/\xi - x/\xi,
\]
n'est pas correctement définie sur $\Xi = [0,1]$; elle signifie aussi que
\[
W =
\begin{cases}
0 & \mbox{ si } \xi = 0; \\
1 & \mbox{ si } \xi \ne 0.
\end{cases}
\]

\end{frame}

\begin{frame}
\frametitle{Exercice 1 - suite}

{\blue (b)} Exprimez $K_2(\xi)$ pour tout $\xi$ dans $[0,1]$.

\mbox{}

Nous allons devoir considérer deux cas: $\xi = 0$ ou $\xi \in (0,1]$.

\mbox{}

{\red 1. $\xi \in (0,1]$} Dans ce cas, comme $y, \xi \geq 0$, $1-x$ doit
être non-négatif pour que le problème soit bien défini:
\[
K_2(\xi) = \lbrace x\,|\, x \leq 1 \rbrace.
\]
La valeur et la solution optimales sont
\[
Q^*(x,\xi) = (1-x)/\xi,\quad y^*= (1-x)/\xi.
\]

{\red 2. $\xi=0$} Il n'existe pas de $y$ tel que $0.y = 1-x$, à moins
que $x = 1$, aussi
\[
K_2(0) = \lbrace 1 \rbrace.
\]

\end{frame}

\begin{frame}
\frametitle{Exercice 1 - suite}

{\blue (c)} Exprimer $K_2$, $K_2^P$ et $\mathcal{Q}$.

\mbox{}

De ce qui précède,
\[
K_2^P = \lbrace x \,|\, x \leq 1 \rbrace \cap \lbrace 1 \rbrace =
\lbrace 1 \rbrace.
\]
Nous avons aussi, vu que $P[\xi = 0] = 0$,
\[
\mathcal{Q}(x) = \int_0^1 \frac{1-x}{\xi} 2\xi d\xi = 2(1-x), \forall
x \leq 1.
\]
Par conséquent $K_2^P \subset K_2 = \lbrace x \leq 1 \rbrace$.

La différence tient au fait qu'un point n'est
pas dans $K_2^P$ dès qu'il est irréalisable pour une certaine valeur
de $\xi$, quelle que soit la distribution de $\bxi$; $K_2$
ne fait pas intervenir les infaisabilités se présentant avec une
probabilité nulle.

\mbox{}

\underline{Rappel}: pour un programme stochastique avec un recours
fixe (pas le cas ici!) où $\bxi$ a des moments seconds finis, $K_2 =
K_2^P$.

\end{frame}

\begin{frame}
\frametitle{Exercice 2}

%(Birge et Louveaux, page 102)
Soit le programme de seconde étape défini par
\begin{align*}
\min_y\ & 2y_1+y_2\\
\mbox{t.q. } & y_1 + y_2 \geq 1-x_1,\\
& y_1 \geq \bxi - x_1-x_2,\\
& y_1, y_2 \geq 0.
\end{align*}
Montrez que ce programme a un recours complet si $\bxi$ est
d'espérance finie.

\mbox{}

Recours complet? Il faut que pour tout $z \in \rit^{2}$, on puisse
trouver $y \geq 0$
t.q. $Wy = z$. Le recours complet est une propriété de $W$, mais vaut
$W$? De manière équivalente, on veut pos$W = \rit^2$.

\end{frame}

\begin{frame}
\frametitle{Exercice 4 - suite}

Petite difficulté: le problème n'est pas sous forme standard.
L'idée est en fait de montrer pour pour tout $z \in \rit^2$, il existe
$y \geq 0$ t.q $Wy \geq z$. L'équivalence avec pos$W$ ne fonctionne
plus. Sauf si nous nous ramenons à la forme standard!

\mbox{}

Sous forme standard, le problème devient
\begin{align*}
\min_y\ & 2y_1+y_2\\
\mbox{t.q. } & y_1 + y_2 - s_1 = 1-x_1,\\
& y_1 - s_2 = \bxi - x_1-x_2,\\
& y_1, y_2 \geq 0,\\
& s_1, s_2 \geq 0,
\end{align*}
et il y a recours complet si pour tout $z \in R^2$, il existe $y$, $s
\geq 0$ vérifiant
\[
Wy = z.
\]

\end{frame}

\begin{frame}
\frametitle{Exercice 2 - suite}

Explicitons $Wy = z$:
\[
\begin{pmatrix}
1 & 1 & -1 & 0 \\
1 & 0 & 0  & -1
\end{pmatrix}
\begin{pmatrix} y_1 \\ y_2 \\ s_1 \\ s_2 \end{pmatrix} =
\begin{pmatrix} z_1 \\ z_2 \end{pmatrix}.
\]
Autrement dit, nous devons trouver $y_1$, $y_2$, $s_1$, $s_2 \geq 0$
tel que
\begin{align*}
y_1 + y_2 - s_1 &= z_1. \\
y_1 - s_2 &= z_2.
\end{align*}

La deuxième équation nous donne la possibilité de prendre
\[
\begin{cases}
y_1 = z_2,\ s_2 = 0 & \mbox{si } z_2 \geq 0,\\
y_1 = 0,\ s_2 = -z_2 & \mbox{sinon}.\\
\end{cases}
\]

\end{frame}

\begin{frame}
\frametitle{Exercice 2 - suite}

Nous avons aussi
\[
y_2 - s_1 = z_1 - y_1.
\]
En utilisant nos précédents choix, nous pouvons prendre pour $y_2$ ce
qui suit. Si $z_2 \geq 0$, $y_2 - s_1 = z_1 - z_2$, et nous prenons
\[
\begin{cases}
y_2 = z_1-z_2,\ s_1 = 0 & \mbox{si } z_1-z_2 \geq 0,\\
y_2 = 0,\ s_2 = z_2-z_1 & \mbox{sinon}.\\
\end{cases}
\]
De même, si $z_2 \leq 0$, $y_2 - s_1 = z_1$, et il suffit de choisir
\[
\begin{cases}
y_2 = z_1,\ s_1 = 0 & \mbox{si } z_1 \geq 0,\\
y_2 = 0,\ s_2 = -z_1 & \mbox{sinon}.\\
\end{cases}
\]

\end{frame}

\begin{frame}
\frametitle{Exercice 2 - suite}

Par conséquent, quel que soit $z$, on peut trouver une solution
non-négative au système
\[
W\begin{pmatrix} y \\ s \end{pmatrix} = z,
\]
i.e. le recours est complet!

%\mbox{}

%Pourquoi exiger $\bxi$ d'espérance finie?

%\mbox{}

%Oups,\ldots pas utile à ce stade (copie de l'erreur dans l'énoncé
%repris de Birge et Louveaux page 102), puisque le recours complet
%n'est une propriété que de $W$, fixe.
La condition d'espérance finie sert surtout à garantir que le problème
de seconde étape est bien défini.
On peut en effet montrer que $\mathcal{Q}(x)$ n'est pas fini si
$\bxi$ n'est pas d'espérance finie (l'expression de $\mathcal{Q}(x)$
fait intervenir $\bxi$).

\end{frame}

\begin{frame}
\frametitle{Exercice 3}

Considérons le programme de deuxième défini comme suit:
\begin{align*}
\min_y\ & 2y_1+y_2\\
\mbox{t.q. } & y_1 - y_2 \leq 2-\bxi x_1,\\
& y_2 \leq x_2,\\
& y_1, y_2 \geq 0.
\end{align*}
Trouvez $K_2(\bxi)$ et $K_2$ pour:
\begin{enumerate}
\item
$\bxi \sim U[0,1]$.
\item
$\bxi \sim \mbox{Poisson}(\lambda),\ \lambda > 0$.
\end{enumerate}
Quelles propriétés attendez-vous pour $K_2$?

\end{frame}

\begin{frame}
\frametitle{Exercice 3 - suite}

Nous devons donc déterminer des ensembles réalisables.

\mbox{}

Remarquons tout d'abord que comme $y_2 \geq 0$ et $y_2 \leq x_2$, nous
devons avoir $x_2 \geq 0$.

\mbox{}

Utilisons la contrainte
\[
y_1 - y_2 \leq 2-\bxi x_1.
\]
Comme $y_2 \leq x_2$, nous avons $y_1 \leq 2-\bxi x_1 + x_2$, et tout
$y_1$ satisfaisant cette inégalité conviendra. Le
problème est donc rélisable si et seulement si $2-\bxi x_1 + x_2 \geq 0$.

Il suit que
\begin{align*}
K_2(\xi) &= \lbrace x = (x_1, x_2)^T \,|\, x_2 \geq 0, \xi x_1 \leq
2+x_2 \rbrace \\
&= \lbrace  x = (x_1, x_2)^T \,|\, x_2 \geq \max(0, \xi x_1 -2 ) \rbrace.
\end{align*}

\end{frame}

\begin{frame}
\frametitle{Exercice 3 - Suite}

Pour les distributions considérées, $E[\xi]$ et $E[\xi^2]$ existent et
sont finies. Il suit que
\[
K_2 = K_2^P = \cap_{\xi \in \Xi} K_2(\xi).
\]

\mbox{}

Si $\bxi \sim U[0,1]$, $\xi \leq 1$, et $\xi x_1 \leq 2 + x_2$ si $x_1
\leq 2 + x_2$. Autrement dit, $K_2(\xi) \subseteq K_2(1)$, et
\[
K_2 = \lbrace x \,|\, x_2 \geq 0, x_1 \leq 2 + x_2 \rbrace.
\]

\mbox{}

Si $\bxi \sim \mbox{Poisson}(\lambda)$, $\xi = 0, 1,\ldots,$, de sorte
que $\xi x_1$ n'est pas borné supérieurement, sauf si $x_1 \leq
0$. Dès lors,
\[
K_2 = \lbrace x \,|\, x_1 \leq 0, x_2 \geq 0 \rbrace.
\]

\end{frame}

\begin{frame}
\frametitle{Exercice 3 - Suite}

Interprétation. Toute décision de première étape doit satisfaire
$x_1$, $x_2 \geq 0$, par conséquent, si $\bxi \sim U[0,1]$,
\[
K_1 \cap K_2 \subseteq \lbrace x \,|\, x_1, x_2 \geq 0, x_1 \leq 2+x_2
\rbrace,
\]
et si $\bxi \sim \mbox{Poisson}(\lambda)$,
\[
K_1 \cap K_2 \subseteq \lbrace x \,|\, x_1 = 0, x_2 \geq 0 \rbrace.
\]
Nous pouvons espérer un recours relativement complet dans le premier
cas (cela dépend de $K_1$), mais le deuxième se présente assez mal!

\end{frame}

\end{document}
